When no active depth sensor is available, Gaussian Splatting SLAM must
derive geometric priors from learned models or visual-inertial
structure-from-motion.  This category---monocular depth + RGB + IMU---is the
most hardware-accessible (requiring only a camera and IMU, both ubiquitous
in consumer devices) but the most algorithmically challenging, as monocular
depth is scale-ambiguous, view-dependent, and noisy.

We subdivide this category into three threads: (1)~visual-inertial
Gaussian SLAM with tightly-coupled IMU fusion, (2)~monocular RGB-only
Gaussian SLAM with learned depth priors, and (3)~feed-forward foundation-model-
based Gaussian SLAM.

\subsection{Visual-Inertial Gaussian Splatting SLAM}

\subsubsection{VIGS-SLAM}

VIGS-SLAM~\cite{zhu2026_vigssslam} (arXiv 2025, ETH Zurich) is the
state-of-the-art visual-inertial 3DGS SLAM system.  It tightly couples
visual and inertial energy terms within a unified optimization framework,
jointly refining camera poses, depths, and IMU states (velocity, bias).
Key contributions include:
\begin{itemize}
    \item \textbf{Custom CUDA kernels} for inertial bundle adjustment that
    directly fuse visual and inertial Hessian contributions in a single
    normal-equation solve per iteration---eliminating the linearization gap
    of cascaded approaches like DBA-Fusion.
    \item \textbf{Robust multi-stage IMU initialization} recovering metric
    scale, gravity direction, and IMU parameters.
    \item \textbf{Time-varying bias modeling} (not fixed bias as in prior
    work).
    \item \textbf{Efficient loop closure} with consistent Gaussian map
    updates.
\end{itemize}
Evaluated on five datasets including a self-captured visual-inertial
dataset, VIGS-SLAM consistently outperforms SOTA methods under motion blur,
low texture, and exposure variations.  \textbf{Project page}:
\url{https://visionlab.ethz.ch/vigs-slam}.

\textbf{Elevator-VIGS}~\cite{zhou2026_elevatorvigs} extends VIGS-SLAM with
a transport state that separates elevator motion from robot motion,
enabling tracking through elevator rides---a scenario where the camera and
IMU observe different motions.

\subsubsection{MM3DGS SLAM}

MM3DGS SLAM~\cite{sun2024_mm3dgs} (IROS 2024) is the first system to use
3D Gaussians for map representation with unposed camera images and inertial
measurements.  Its keyframe-based framework incorporates relative pose
transformations from pre-integrated IMU measurements, depth estimates, and
photometric rendering quality.  For monocular input, it treats learned
monocular depth as sensor depth.  Achieves 3$\times$ tracking improvement
over prior 3DGS SLAM SOTA.

\subsubsection{CGS-SLAM}

CGS-SLAM~\cite{deambrogi2026_cgssslam} (arXiv 2026) enables multi-agent
3DGS SLAM using only RGB and inertial data.  Each agent performs local
tracking with inertial data as a motion prior and reconstructs a scaled map
using a \emph{metric} monocular depth estimator (Depth Pro).  Keyframe
encodings are shared among agents for dynamic keyframing in overlap
regions, and a central server aligns submaps using VGGT~\cite{wang2025_vggt}
as a view alignment model.  This is the collaborative counterpart to
VIGS-SLAM, targeting GNSS-denied tactical scenarios with consumer
smartphones.

\subsubsection{VIM-GS}

VIM-GS~\cite{zhang2026_vimgs} (arXiv 2025) is a visual-inertial monocular
Gaussian Splatting framework for large scenes.  It addresses the depth
challenge by fusing accurate but sparse depth from visual-inertial SfM with
dense but coarse depth from large foundation models (LFMs).  An
object-segmented depth propagation algorithm bridges sparse and dense
depth, and a dynamic depth refinement module handles crippled SfM depth of
dynamic objects.  This is the large-scale outdoor counterpart to indoor
VIGS-SLAM.

\subsection{Monocular RGB-Only Gaussian SLAM with Learned Depth}

\subsubsection{MonoGS (Gaussian Splatting SLAM)}

MonoGS~\cite{matsuki2024monogs} (CVPR 2024) is the first application of 3D
Gaussian Splatting in monocular SLAM.  It uses 3DGS as the sole scene
representation and introduces geometric regularization to address
ambiguities in incremental reconstruction from monocular input.  The system
demonstrates near real-time operation and achieves SOTA results in novel
view synthesis and trajectory estimation, including reconstruction of tiny
and transparent objects.  \textbf{Code}:
\url{https://github.com/muskie82/MonoGS}.

\subsubsection{HI-SLAM2}

HI-SLAM2~\cite{zhang2026_hislam2} (IEEE T-RO 2025) achieves fast and
accurate monocular scene reconstruction using only RGB input.  Its key idea
is combining easy-to-obtain monocular geometry priors (depth + normals from
OmniData) with learning-based dense SLAM, using 3DGS as the map
representation.  A grid-based scale alignment strategy addresses scale
distortions in monocular depth priors.  Upon loop closure, efficient
pose-graph BA with instant Gaussian deformation ensures global consistency.
Surpasses even RGB-D methods on Replica, ScanNet, and ScanNet++.  \textbf{Code}:
\url{https://hi-slam2.github.io/}.

\subsubsection{Splat-SLAM}

Splat-SLAM~\cite{sandstrom2025_splatslam} (CVPR 2025) is the first RGB-only
SLAM system with globally optimized tracking, dynamically adapting the
Gaussian map to keyframe pose and depth updates.  It incorporates
Disparity, Scale and Pose Optimization (DSPO) for bundle adjustment,
jointly optimizing pose, depth, and monocular depth scale.

\subsubsection{GaussianFlow SLAM}

GaussianFlow SLAM~\cite{seo2026_gaussianflow} (IEEE RA-L 2026) leverages
optical flow as a geometry-aware cue to guide the optimization of both the
3DGS and camera poses.  By aligning the projected motion of Gaussians
(``GaussianFlow'') with dense optical flow, it continuously corrects 3DGS
geometry even without direct depth measurements.  Normalized error-based
densification and pruning modules refine inactive Gaussians.  \textbf{Code}:
\url{https://github.com/url-kaist/gaussianflow-slam}.

\subsubsection{MyGO-Splat}

MyGO-Splat~\cite{zhu2026_mygosplat} (arXiv 2026) is a closed-loop Gaussian
SLAM framework that analytically rasterizes Gaussian primitives into
pixel-wise depth and surface normals, allowing the map to actively supervise
camera pose optimization.  Achieves performance comparable to RGB-D methods
while using only monocular input.

\subsubsection{WildGS-SLAM}

WildGS-SLAM~\cite{zheng2025_wildgssslam} (CVPR 2025) addresses monocular
Gaussian SLAM in dynamic environments.  An uncertainty map, predicted by a
shallow MLP given DINOv2 features, guides dynamic object removal during
tracking (weighted DBA) and mapping (weighted rendering loss).  Monocular
metric depth facilitates pose estimation.  A new Wild-SLAM dataset with
diverse indoor/outdoor dynamic scenes is contributed.

\subsubsection{GeoGS-SLAM}

GeoGS-SLAM~\cite{zhou2026_geogs} (arXiv 2026) introduces geometry-only
Gaussian Splatting (GeoGS) that directly reconstructs scene geometry, and
builds a dense visual SLAM system upon it.  It outperforms SOTA methods in
online mapping efficiency and geometric reconstruction quality.

\subsubsection{MonoEM-GS}

MonoEM-GS~\cite{kruzhkov2026_monoemgs} (arXiv 2026) integrates feed-forward
geometric foundation model predictions into a global Gaussian Splatting
representation via an Expectation-Maximization formulation that explicitly
addresses view-dependent geometry inconsistencies.  ICP-based alignment
provides monocular pose estimation.  Parameterizes Gaussians with
multi-modal features for in-place open-set segmentation.

\subsection{Feed-Forward Foundation-Model-Based Gaussian SLAM}

\subsubsection{VGGT-GS SLAM}

VGGT-GS SLAM~\cite{han2026_vggtgssslam} (arXiv 2026) is a monocular 3DGS
SLAM system for uncalibrated videos.  Starting from feed-forward VGGT
pose and depth priors, it performs submap differentiable bundle adjustment
that jointly refines camera poses, a 3D Gaussian map, and
submap-shared intrinsics and radial-tangential distortion through analytic
calibration Jacobians.  Gaussian-native alignment (GNA) provides
camera-anchored scale refinement between submaps and loop-closure
verification.

\subsection{Summary of Monocular-Depth-Based Systems}

\Cref{tab:mono_summary} compares the key monocular-depth Gaussian SLAM
systems.  The field has progressed rapidly from the initial MonoGS
(RGB-only, near real-time) to VIGS-SLAM (tightly-coupled VI, real-time, loop
closure) and foundation-model-based systems (VGGT-GS SLAM, uncalibrated).

\begin{table}[t]
\centering
\caption{Monocular Depth + RGB + IMU Gaussian Splatting SLAM systems.}
\label{tab:mono_summary}
\scriptsize
\begin{tabular}{lcccccc}
\toprule
\textbf{System} & \textbf{Year} & \textbf{Venue} & \textbf{IMU} & \textbf{Depth Prior} & \textbf{Loop} & \textbf{Code} \\
\midrule
MonoGS~\cite{matsuki2024monogs} & 2024 & CVPR & -- & None & \checkmark & \checkmark \\
MM3DGS~\cite{sun2024_mm3dgs} & 2024 & IROS & \checkmark & Mono depth & -- & -- \\
HI-SLAM2~\cite{zhang2026_hislam2} & 2025 & T-RO & -- & OmniData & \checkmark & \checkmark \\
Splat-SLAM~\cite{sandstrom2025_splatslam} & 2025 & CVPR & -- & Mono depth & \checkmark & -- \\
WildGS-SLAM~\cite{zheng2025_wildgssslam} & 2025 & CVPR & -- & Metric depth & -- & -- \\
VIM-GS~\cite{zhang2026_vimgs} & 2025 & arXiv & \checkmark & LFM + SfM & -- & -- \\
VIGS-SLAM~\cite{zhu2026_vigssslam} & 2025 & arXiv & \checkmark & DROID-SLAM & \checkmark & -- \\
GaussianFlow~\cite{seo2026_gaussianflow} & 2026 & RA-L & -- & Optical flow & -- & \checkmark \\
GeoGS-SLAM~\cite{zhou2026_geogs} & 2026 & arXiv & -- & Learned & -- & -- \\
MyGO-Splat~\cite{zhu2026_mygosplat} & 2026 & arXiv & -- & Analytic & -- & -- \\
MonoEM-GS~\cite{kruzhkov2026_monoemgs} & 2026 & arXiv & -- & Foundation & -- & -- \\
VGGT-GS SLAM~\cite{han2026_vggtgssslam} & 2026 & arXiv & -- & VGGT & \checkmark & -- \\
CGS-SLAM~\cite{deambrogi2026_cgssslam} & 2026 & arXiv & \checkmark & Depth Pro & -- & -- \\
\bottomrule
\end{tabular}
\end{table}

\textbf{Key insight}: The fundamental challenge in monocular-depth Gaussian
SLAM is \emph{scale ambiguity}.  Three approaches address this:
(a)~IMU pre-integration for metric scale (VIGS-SLAM, MM3DGS, CGS-SLAM),
(b)~grid-based scale alignment of monocular depth priors (HI-SLAM2), and
(c)~metric monocular depth estimation models (Depth Pro in CGS-SLAM,
Metric3D in MAGS-SLAM).  The IMU-based approach is most robust but requires
hardware; the metric depth model approach is hardware-free but relies on
generalization of the depth network.
